% Answer key for the Module 2 pen-and-paper sheet.
% Compile with: xelatex -interaction=nonstopmode solutions.tex   (run twice)
%
% It carries the questions as well as the answers, so it can be read on its
% own, and it draws the same town as the sheet out of townkit.tex -- an answer
% to a question about a drawing should be a drawing, and fifteen handshake
% counts in a row of blanks is the one thing that cannot show a student the
% shape they were meant to see.
\documentclass[a4paper, 11pt]{extarticle}
\usepackage[top=0.7in, bottom=0.7in, left=0.8in, right=0.8in]{geometry}
\usepackage{amsmath}
\usepackage{amssymb}
\usepackage{array}
\usepackage{graphicx}
\usepackage{tcolorbox}
\usepackage{enumitem}
\usepackage{etoolbox}
\usepackage{tikz}
\usepackage{fontspec}
\usetikzlibrary{calc,positioning}
\definecolor{pltblue}{HTML}{1F77B4}
\definecolor{pltrust}{HTML}{B14434}

% Charter, and the four ways to get it. XCharter *is* Charter and ships with a
% full TeX Live, so it is asked for first and by file: a CI runner answers yes
% to \IfFontExistsTF{Charter}, then cannot typeset it, which is how this sheet
% stopped building there.
\IfFontExistsTF{XCharter-Roman.otf}{
  \setmainfont{XCharter}[
    Extension      = .otf,
    UprightFont    = *-Roman,
    BoldFont       = *-Bold,
    ItalicFont     = *-Italic,
    BoldItalicFont = *-BoldItalic]
}{
  \IfFontExistsTF{Charter}{\setmainfont{Charter}}{
    \IfFontExistsTF{Iowan Old Style}{\setmainfont{Iowan Old Style}}{
      \IfFontExistsTF{Georgia}{\setmainfont{Georgia}}{}}}}

\newcommand{\ppfont}{}
\IfFontExistsTF{Pretty Neat}{\renewcommand{\ppfont}{\fontspec{Pretty Neat}}}{
  \IfFontExistsTF{Humor Sans}{\renewcommand{\ppfont}{\fontspec{Humor Sans}}}{
    \IfFontExistsTF{Excalifont}{\renewcommand{\ppfont}{\fontspec{Excalifont}}}{
      \renewcommand{\ppfont}{}}}}

\setlength{\parindent}{0pt}
\setlength{\parskip}{0.4em}
\raggedbottom

\newcommand{\parthead}[2]{%
  \vspace{0.5em}
  {\large\bfseries Part #1 --- #2}\par
  \vspace{-0.4em}\rule{\textwidth}{0.8pt}\par
  \vspace{0.2em}}

\newcolumntype{C}[1]{>{\centering\arraybackslash}p{#1}}

% An answer, set apart from the question it belongs to. Wrapped rather than
% switched: a bare \color inside a p-column cell costs the cell its first line.
\newcommand{\ans}[1]{\textbf{\textcolor{pltblue!75!black}{#1}}}

% The line that turns from question to answer.
\newcommand{\arrow}{\textcolor{pltblue!75!black}{$\Rightarrow$}\;}

% The same as \ans, for something already inside math. \ans wraps its argument
% in \textbf, which leaves math mode and takes the formula with it.
\newcommand{\mans}[1]{{\color{pltblue!75!black}#1}}

% Ringville and the neighbourhood figure, shared with exercise.tex.
% Ringville, and the two other figures the Module 2 sheet draws. Kept in their
% own file because the answer key draws exactly the same town with the boxes
% filled in. Needs tikz, its calc and positioning libraries, and the colours
% pltblue and pltrust, all of which the including file sets up.
% ---------------------------------------------------------------------------
% The town: sixteen people in a circle, everybody friends with the two on
% their left and the two on their right. Four friends each, thirty-two
% friendships. That one drawing has to carry both halves of the module --
% distance AND triangles -- which is why it is next-two-neighbours rather than
% next-one: a ring of next-one neighbours has no triangle anywhere in it.
%
% Two more things fall out of the choice and both are used on the sheet. The
% town is left-right symmetric about the line through person 1 and person 9,
% so a student who has worked out people 2 to 9 already has 10 to 16. And the
% two shortcuts (1--9 and 5--13) are diameters, so they cross at the middle of
% the drawing and read as two lines rather than as a tangle.
% ---------------------------------------------------------------------------
\tikzset{
  % Every friendship is drawn the same, so no line weight says "this one
  % matters more". The two shortcuts are the exception, and they are the point.
  friend/.style   = {line width=1.4pt, black!58, line cap=round},
  shortcut/.style = {line width=2.1pt, pltrust, line cap=round},
  chair/.style    = {circle, draw=black!72, fill=white, line width=0.9pt,
                     minimum size=0.66cm, inner sep=0pt, font=\footnotesize},
  chairhot/.style = {chair, draw=pltblue!75!black, line width=1.5pt,
                     fill=pltblue!14},
  % The box a student writes a handshake count into, one per person, printed
  % outside the ring where no friendship line can reach it.
  dbox/.style     = {draw=pltblue!70, line width=0.7pt, fill=white,
                     minimum size=0.58cm, inner sep=0pt, rounded corners=1.2pt},
  dnum/.style     = {font=\bfseries\footnotesize, inner sep=0pt},
  % A pair that might or might not be a friendship. The student darkens it.
  maybe/.style    = {line width=1.3pt, black!52, dash pattern=on 1.6pt off 2.2pt},
}

% \ringtown{extra lines}{numbers already written in}.
%
% The count goes in the box beside the person, on the drawing, rather than
% into a list of blanks underneath it: the thing the sheet is teaching is that
% the number belongs to a seat, and a column of fifteen blanks hides that.
\newcommand{\ringtown}[2]{%
\begin{tikzpicture}[x=1cm, y=1cm]
    \foreach \i in {1,...,16} {
        \pgfmathsetmacro{\ang}{90 - (\i-1)*22.5}
        \coordinate (c\i) at (\ang:2.7);
    }
    % Next door, both sides: the rim of the circle.
    \foreach \i in {1,...,16} {
        \pgfmathtruncatemacro{\j}{mod(\i,16)+1}
        \draw[friend] (c\i) -- (c\j);
    }
    % One seat over, both sides. Bowed inward: drawn straight, the chord passes
    % 0.14cm inside the chair sitting between its two ends and is drawn through
    % that person's face.
    \foreach \i in {1,...,16} {
        \pgfmathtruncatemacro{\j}{mod(\i+1,16)+1}
        \pgfmathsetmacro{\amid}{90 - \i*22.5}
        \draw[friend] (c\i) .. controls (\amid:2.12) .. (c\j);
    }
    #1
    \foreach \i in {1,...,16} {
        \pgfmathsetmacro{\ang}{90 - (\i-1)*22.5}
        \node[chair] (p\i) at (\ang:2.7) {\i};
        \node[dbox, name=b\i] at (\ang:3.44) {};
    }
    #2
\end{tikzpicture}}

% The two friendships that came back from college. Both are diameters of the
% circle -- 1 is opposite 9 and 5 is opposite 13 -- so they cross at the
% centre, which is the one place in the drawing with nothing else in it.
\newcommand{\shortcuts}{%
    \draw[shortcut] (c1) -- (c9);
    \draw[shortcut] (c5) -- (c13);}

% A number written into person #1's box.
\newcommand{\dfill}[2]{\node[dnum] at (b#1) {#2};}

% Every distance from person 1 in the plain town, and in the town with the two
% shortcuts. The sheet prints two of each; the answer key prints them all.
\newcommand{\plainfill}{%
  \dfill{1}{0}\dfill{2}{1}\dfill{3}{1}\dfill{4}{2}\dfill{5}{2}\dfill{6}{3}%
  \dfill{7}{3}\dfill{8}{4}\dfill{9}{4}\dfill{10}{4}\dfill{11}{3}\dfill{12}{3}%
  \dfill{13}{2}\dfill{14}{2}\dfill{15}{1}\dfill{16}{1}}
\newcommand{\cutfill}{%
  \dfill{1}{0}\dfill{2}{1}\dfill{3}{1}\dfill{4}{2}\dfill{5}{2}\dfill{6}{3}%
  \dfill{7}{2}\dfill{8}{2}\dfill{9}{1}\dfill{10}{2}\dfill{11}{2}\dfill{12}{3}%
  \dfill{13}{2}\dfill{14}{2}\dfill{15}{1}\dfill{16}{1}}

% ---------------------------------------------------------------------------
% Person 3's four friends, lifted out of the circle. Person 3 is NOT in this
% picture: drawn in, their four lines cross the six pairs the student is being
% asked to read, and the figure stops being readable at exactly the moment it
% is needed.
%
% The four are placed 1, 4, 2, 5 clockwise from the top left rather than in
% seat order. In seat order the three real friendships are the three shortest
% lines in the figure and the answer can be had without looking at the town;
% shuffled, the student has to go back to the circle for all six.
% ---------------------------------------------------------------------------
\newcommand{\hoodfig}[1]{%
\begin{tikzpicture}[x=1cm, y=1cm]
    \coordinate (h1) at (-1.55, 1.30);
    \coordinate (h4) at ( 1.55, 1.30);
    \coordinate (h2) at ( 1.55,-1.30);
    \coordinate (h5) at (-1.55,-1.30);
    \draw[maybe] (h1) -- (h4);
    \draw[maybe] (h4) -- (h2);
    \draw[maybe] (h2) -- (h5);
    \draw[maybe] (h5) -- (h1);
    \draw[maybe] (h1) -- (h2);
    \draw[maybe] (h4) -- (h5);
    #1
    \foreach \i in {1,2,4,5} {\node[chair] at (h\i) {\i};}
\end{tikzpicture}}

% The three of the six that really are friendships, for the answer key.
\newcommand{\hoodreal}{%
    \draw[friend, pltblue!80!black] (h4) -- (h2);
    \draw[friend, pltblue!80!black] (h1) -- (h2);
    \draw[friend, pltblue!80!black] (h4) -- (h5);}


\begin{document}

\begin{center}
{\ppfont\LARGE Answers}
\end{center}

{\small The sheet with its answers, \ans{in blue}. Question 1 is a guess and
has no wrong answer; every other question has one right one, except 8(c),
where the chain depends on a tie you are free to break either way.}

\parthead{1}{A letter from Omaha}

{\bf Question 1}: The letter arrived. How many hands did it pass through on the
way? Write your guess, and write why you guessed that.

\arrow A guess, and the room's guesses are the material for the session. The
number is in 8(d): the chains that arrived ran about \ans{five or six} people
long. Most guesses in a lecture hall are far too high, which is exactly what
makes the arithmetic in 8(a) worth doing.

\parthead{2}{Ringville}

Sixteen people live in Ringville and they sit in one circle. Everybody is
friends with the two on their left and the two on their right: four friends
each, thirty-two friendships in the town.

{\bf Question 2(a)}: Person 1 has a letter for person 9. Trace the chain on the
drawing. How many handshakes does it take?

\arrow \ans{Four}, and there are exactly two chains that do it:
\ans{1--3--5--7--9} and \ans{1--15--13--11--9}. Person 9 is eight chairs away
and no single handshake covers more than two chairs, so nothing beats four.

{\bf Question 2(b)}: Now everybody. In each person's box on the drawing, write
how many handshakes it takes to reach them from person 1. Two are written in
already, and you may find you only have to work out half of the rest.

\arrow Below left. Going round one way gives the same numbers as going round
the other, so people 10 to 16 are people 8 to 2 read backwards --- half the
work, and the reason the drawing is symmetric is that person 1 has no idea
which way is which.

\begin{center}
\begin{tabular}{@{}c@{\hspace{0.9cm}}c@{}}
\resizebox{0.43\textwidth}{!}{\ringtown{}{\plainfill}} &
\resizebox{0.43\textwidth}{!}{\ringtown{\shortcuts}{\cutfill}} \\[0.4em]
\small\bfseries 2(b): Ringville & \small\bfseries 4: with the two shortcuts \\
\end{tabular}
\end{center}

{\bf Question 2(c)}: Add up every number in the boxes. What is the total?
Share that total out among the fifteen other people --- the town's average. And
what is the biggest number in any box, the town's worst case?

\arrow Total \ans{36}, average $36/15 = \ans{2.4}$, worst case \ans{4}, which
three people share: 8, 9 and 10.

{\bf Question 2(d)}: Take the average and the worst case. One of the two you
would give to a newspaper; the other you would give to an engineer who has to
promise the letter arrives. Which is which, and why?

\arrow \ans{2.4 is the newspaper's} number and \ans{4 is the engineer's}. An
average describes the letter you would typically send; it promises nothing
about any particular one. The worst case is the only number you can put a
guarantee behind, because no letter in this town ever takes longer. Both get
names in the lecture: \ans{average path length} and \ans{diameter}.

{\bf Question 3}: Ringville grows. A thousand people, still one circle, still
four friends each. How many handshakes does it take to reach the person sitting
directly opposite you? And if the circle held all eight billion people on Earth?
Could a world wired like Ringville get the Omaha letter to Boston in six
handshakes?

\arrow With a thousand people the person opposite is 500 chairs away, so
\ans{250} handshakes. With eight billion they are four billion chairs away, so
\ans{two billion}. \ans{No} --- in a ring the distance grows in step with the
town, so a ring-wired world would need billions of handshakes, not six.
Whatever explains Milgram, it is not this shape.

\parthead{3}{Two shortcuts}

Two of the sixteen went to college and came home with a friend from the far
side of town: person 1 is now also friends with person 9, and person 5 with
person 13 --- thirty-four friendships, not thirty-two.

{\bf Question 4}: Write the handshake counts into the boxes again, this time for
the town with the two new friendships. Then the same three numbers as in 2(c):
the new total, the new average, and the new worst case.

\arrow Above right. Total \ans{27}, average $27/15 = \ans{1.8}$, worst case
\ans{3}, now reached only by persons 6 and 12.

{\bf Question 5}: Thirty-two friendships became thirty-four. Work out both
changes as percentages: by how much did the number of friendships change, and by
how much did the average change? Say what those two new friendships did that
thirty-two ordinary ones could not.

\arrow Friendships up by $2/32 = \ans{6.25\%}$; the average down from 2.4 to
1.8, which is \ans{25\%}. Four times the effect, for one sixteenth of the
wiring.

An ordinary Ringville friendship joins two people who were already one or two
handshakes apart, so it shortens nothing for anybody. These two joined people
\emph{four} handshakes apart --- and the saving is not theirs alone. Once
person 1 can reach person 9 in a single step, so can everybody sitting near
person 1, for everybody sitting near person 9. One line rewrote the distance
between two whole neighbourhoods.

\parthead{4}{Do your friends know each other?}

Go back to Ringville as it was in Part 2, before the two shortcuts, and take
somebody ordinary. Person 3's four friends are persons 1, 2, 4 and 5, and those
four make six pairs.

{\bf Question 6(a)}: Check the six pairs against the circle in Part 2, one pair
at a time. Darken the dashed lines that are real friendships. How many of the
six are?

\arrow \ans{Three of six}, drawn below: \ans{1 \& 2}, \ans{2 \& 4} and
\ans{4 \& 5}. The other three --- 1 \& 4, 1 \& 5, 2 \& 5 --- sit three or four
chairs apart, which is further than a Ringville friendship reaches.

\begin{center}
\resizebox{0.26\textwidth}{!}{\hoodfig{\hoodreal}}
\end{center}

{\bf Question 6(b)}: Ringville grows to ten thousand people, still four friends
each. How many of person 3's six pairs are friends now?

\arrow \ans{Still three of six.} Person 3's friends are still the people in
chairs 1, 2, 4 and 5, and whether two of them are friends depends only on how
many chairs apart they sit. The size of the town never enters the question.

{\bf Question 6(c)}: A different town: the same sixteen people and the same
thirty-two friendships, but this time drawn out of a hat. Pick two people. What
is the chance that they are friends? So how many of person 3's six pairs would
you expect to be friends? And in a hat-drawn town of ten thousand, still four
friends each?

\arrow Each person has 4 friends among the 15 others, so the chance is
$4/15 \approx \ans{0.27}$, and $6 \times 4/15 = \ans{1.6}$ of the six pairs.
In a hat-drawn town of ten thousand it is $4/9999$, about one in 2{,}500, so
$6 \times 4/9999 = \ans{0.0024}$ --- \ans{effectively none}.

{\bf Question 6(d)}: That fraction has a name --- it is person 3's
\textbf{local clustering}. You have now measured it four times. Say in one
sentence what makes it big and what makes it small.

\arrow It is \ans{big when your friends are drawn from the same small pocket of
the world you are drawn from}, because then they are in each other's pockets
too, and \ans{small when they are drawn from the whole town}. Ringville holds
$1/2$ at any size; the hat-drawn town falls away as $1/n$. That is the whole
reason the two halves of this module are a puzzle at all: \ans{short paths come
free from randomness, and high clustering does not}.

Three people who are all friends with each other make a \textbf{triangle}.
Ringville, before the two shortcuts, has sixteen.

{\bf Question 7(a)}: Then the two shortcuts went in. Do persons 1 and 9 have a
friend in common? Do persons 5 and 13? So how many triangles does the town have
now?

\arrow Person 1's friends are 15, 16, 2, 3; person 9's are 7, 8, 10, 11 ---
\ans{no overlap}, and \ans{none} for 5 and 13 either, so neither new line
closes a triangle. And a line that is \emph{added} can never destroy one. So
\ans{sixteen, unchanged}. The average fell by a quarter and the town lost
nothing.

{\bf Question 7(b)}: The real model does not \emph{add} a shortcut, it
\textbf{moves} one: person 1 drops person 2 and befriends person 9 instead.
Which friends did persons 1 and 2 have in common? So how many triangles does
the town lose?

\arrow Persons 1 and 2 both know \ans{person 16 and person 3}, so cutting
1--2 kills the triangles \ans{1--2--16} and \ans{1--2--3}: \ans{two} go, and
the town is left with \ans{fourteen}. This is what the Watts--Strogatz model
actually does at each step, and it is the honest version of 7(a).

{\bf Question 7(c)}: Put Question 5 beside 7(a) and 7(b). What does a shortcut
buy the town, and what does it cost?

\arrow It buys \ans{a quarter off the average distance} for two lines out of
thirty-four. It costs \ans{nothing when the line is added, and two triangles
out of sixteen when it is moved} --- and the triangles it breaks are all local
to the one person it was moved from. \ans{Distance is cheap to buy and
clustering is expensive to lose}, so a handful of shortcuts collapses the
distances and leaves almost all of the local structure standing. That is the
mechanism the whole model rests on.

\parthead{5}{Back to Omaha}

{\bf Question 8(a)}: Everybody knows 150 people. Pretend for a moment that no
two of your friends know each other, so every handshake reaches 150 people
nobody in the chain has met yet. Carry the table two rows further. At how many
handshakes do you pass eight billion?

\begin{center}
{\setlength{\tabcolsep}{9pt}
\begin{tabular}{|C{2.6cm}|C{4.6cm}|}
\hline
{\small\bfseries handshakes} & {\small\bfseries people reached} \\
\hline
1 & 150 \\ \hline
2 & 22,500 \\ \hline
3 & about 3.4 million \\ \hline
4 & \ans{about 506 million} \\ \hline
5 & \ans{about 76 billion} \\ \hline
\end{tabular}}
\end{center}

\arrow \ans{Five.} Four handshakes reach half a billion, which is not enough;
five overshoot the planet by a factor of nine. Six degrees was never a
measurement, it was this multiplication.

{\bf Question 8(b)}: In Question 6 you found that person 3's friends do know
each other, so the pretending in 8(a) is false. Which way does that push the
answer --- more handshakes, or fewer? Why?

\arrow \ans{More handshakes.} If your friends know each other --- and Question
6 says half their pairs do --- then most of the 150 people your friend reaches
are people you were already reaching yourself. Every step adds far fewer than
150 times as many new people, so the tower of powers above is an
\ans{overestimate of the reach} and therefore an \ans{underestimate of the
number of handshakes}. Clustering is exactly what makes six degrees hard.

{\bf Question 8(c)}: You are person 3 in the town with the two shortcuts,
holding a letter for person 12. You know only your own four friends and the
chair each of them sits in, and you hand it to whichever of them sits closest to
chair 12; they do the same. Write the chain that rule produces, then the
shortest chain that exists.

\arrow Person 3's friends sit in chairs 1, 2, 4 and 5, which are 5, 6, 8 and 7
chairs from chair 12 --- so the letter goes to \ans{person 1}. Person 1's
friends are 15, 16, 2, 3 and 9, and \ans{15 and 9 tie} at three chairs away.
Either way it takes four:
\begin{center}
\ans{3 --- 1 --- 15 --- 13 --- 12} \qquad or \qquad \ans{3 --- 1 --- 9 --- 11
--- 12} \qquad \ans{4 handshakes}
\end{center}
The shortest chain is \ans{3 --- 5 --- 13 --- 12}, \ans{3 handshakes}.

Person 3 never sends it that way, and is not being stupid. Person 5 sits
\emph{seven} chairs from the target --- further round than person 1 --- so by
the only rule person 3 can apply, person 5 is the wrong way. What person 3
cannot know is that person 5 is the one holding the door to 13.

{\bf Question 8(d)}: Of 296 letters sent from Omaha, 64 arrived, in five or six
steps. Is ``a short chain exists'' the same claim as ``an ordinary person can
find one''? Which of the two did the letters test, and what do the 232 that
never arrived say about it?

\arrow \ans{Not the same claim}, and 8(c) is a town where the second is false
while the first is true. The letters tested \ans{the finding, not the
existing}: nobody in the chain ever saw the network, and each of them made one
local guess.

And \ans{232 of the 296 never arrived at all}. A chain that stalls is not
evidence that no short chain existed --- it is evidence that somebody could not
see it, or could not be bothered. The experiment measured the length of the
chains that \emph{worked}, which is a number about people, not only about the
network. That distinction is what the lecture is for.

\clearpage

\parthead{6}{The lab}

Three functions, and everything else in the notebook runs off them.

\begin{verbatim}
def ring_edges(n, half):
    return [(i, (i + d) % n) for i in range(n) for d in range(1, half + 1)]

def distances_from(A, s):
    n = len(A)
    dist = np.full(n, -1)
    dist[s] = 0
    frontier = [s]
    while frontier:
        nxt = []
        for u in frontier:
            for v in np.flatnonzero(A[u]):
                if dist[v] < 0:
                    dist[v] = dist[u] + 1
                    nxt.append(v)
        frontier = nxt
    return dist

def local_clustering(A, i):
    nbrs = np.flatnonzero(A[i])
    k = len(nbrs)
    if k < 2:
        return 0.0
    links = A[np.ix_(nbrs, nbrs)].sum() / 2
    return links / (k * (k - 1) / 2)
\end{verbatim}

Run on the sheet's town they give \ans{36 / 2.4 / 4} and \ans{0.5} for every
person, which are Questions 2(c) and 6(a) again, out of code that never looked
at the drawing.

\vspace{0.4em}

{\bf Finished early --- the number that lies.} The notebook then measures
$\sigma = (C/C_{\text{rand}}) / (L/L_{\text{rand}})$ on a plain ring of
\emph{one thousand} people with \ans{no shortcuts at all}, and gets
\ans{$\sigma = 4.96$}: a strong small world, by a test applied to a town that
has none. Turn the $n$ dial up and the verdict gets \emph{stronger}, not
weaker.

The four quantities, for a ring of $n$ people with $k$ friends each:
\[
C = \tfrac12, \qquad L \approx \tfrac{n}{8}, \qquad
C_{\text{rand}} = \tfrac{k}{n}, \qquad
L_{\text{rand}} \approx \tfrac{\ln n}{\ln k}
\]
\[
\frac{C}{C_{\text{rand}}} = \frac{n}{8}, \qquad
\frac{L}{L_{\text{rand}}} = \frac{n}{8}\cdot\frac{\ln k}{\ln n}, \qquad
\sigma = \frac{n/8}{(n/8)(\ln k/\ln n)} = \mans{\frac{\ln n}{\ln k}}
\]

The $n$ cancels exactly. A ring is $n/8$ times more clustered than a random
town \emph{and} $n/8$ times further across, and $\sigma$ divides one by the
other, so it has measured nothing about this town at all. What survives the
cancellation is the leftover $\ln n / \ln k$, which is there only because a
random town's paths grow like $\ln n$ --- so slowly that the ring's enormous
path-length penalty is discounted by that factor, and the clustering side wins
by it. Any network whose clustering does not fall with $n$ scores
$\sigma > 1$ automatically, shortcuts or no shortcuts. \ans{$\sigma$ was
certifying that $C$ is constant, not that shortcuts exist.}

The published repairs both add a \emph{second} yardstick --- a lattice, at the
other end from the random graph --- because ``small world'' was always a claim
about the middle of a range, and $\sigma$ only ever looked at one end:

\begin{center}
$\omega = L_{\text{rand}}/L - C/C_{\text{latt}}$ \quad (Telesford 2011;
$-1$ lattice, $0$ small world, $+1$ random)
\end{center}

At $n = 1000$, $k = 4$:

\begin{center}
{\setlength{\tabcolsep}{10pt}
\begin{tabular}{|l|C{2.4cm}|C{2.4cm}|C{2.4cm}|}
\hline
\textbf{town} & $\boldsymbol\sigma$ & $\boldsymbol\omega$ & \textbf{SWI} \\
\hline
plain ring, $p=0$ & \ans{4.96} \ passes & $-0.96$ & 0.00 \\ \hline
rewired, $p=0.05$ & 47.7 & $-0.41$ & \ans{0.81} \ largest \\ \hline
rewired, $p=0.2$  & 48.1 & 0.21 & 0.51 \\ \hline
random, $p=1$     & 0.47 & 0.93 & 0.00 \\ \hline
\end{tabular}}
\end{center}

The same network is a \emph{strong small world} to $\sigma$ and an
\emph{almost perfect lattice} to $\omega$, and as $n$ grows the two run in
opposite directions. The price of the repair is that $C_{\text{latt}}$ and
$L_{\text{latt}}$ have to come from somewhere: here the town \emph{is} a
lattice so they are free, and on real data you have to build a lattice with
the same degrees yourself, which is not a small job. That is part of why
$\sigma$ is still the one people quote.

\vspace{0.4em}

{\small The whole notebook with every blank filled in is
\texttt{lab-solutions.py}, built from \texttt{lab.py} by\\
\texttt{tools/build\_lab\_notebooks.py m02-small-world}.}

\end{document}
